\subsection{TD-MPC2: separating proposal and scoring consumers}
\label{sec:tdmpc-consumer-factorial}
We test a different explanation for an existing development-case return gap:
can a dynamics update cause harm primarily by changing the policy proposals,
so that restoring the scoring path cannot recover performance?
We reuse two saved dynamics modules fitted from the random-source ($R$) and
cascade-source ($C$) \texttt{cheetah-run} trajectories at source seed 23.
$R$ is an adapted comparison model, neither the original pretrained Clean
model nor an oracle for physical outcomes. No new fitting is performed.

The native TD-MPC2 planner uses dynamics both to generate 24 policy-prior
trajectories and to evaluate its 512 candidates. We independently assign
$R$ or $C$ to these two consumers, preserving the native horizon 3,
six optimization iterations, 64 elites, stochastic action selection,
normalization, and all other original model weights.
The scoring route includes imagined-state reward readouts and terminal
policy/$Q$ evaluation; it is not a direct forecast-distance score.
Each arm executes one complete 500-control \texttt{walker-run} episode from
seed 31, with the same initial state and reset random streams.
Later states, proposal populations, and optimizer warm starts may diverge
as endogenous consequences of the intervention.

\begin{table}[ht]
\centering
\begin{tabular}{lccr}
\toprule
Arm & Proposal dynamics & Scoring dynamics & Actual episode return \\
\midrule
$RR$ & $R$ & $R$ & 612.082555 \\
$CC$ & $C$ & $C$ & 495.216321 \\
$CR$ & $C$ & $R$ & 674.885888 \\
$RC$ & $R$ & $C$ & 553.538964 \\
\bottomrule
\end{tabular}
\caption{One reused TD-MPC2 development case. Returns are sums of actual
500-control environment rewards, not planner value estimates.}
\label{tab:tdmpc-consumer-factorial}
\end{table}

Both $RR$ and $CC$ exactly reproduce the saved observations, actions,
per-control rewards, and complete returns. At the common first query,
all arms share the 24 prior first actions and all 488 Gaussian candidate
trajectories exactly; prior continuation actions can differ. Recording these
already-generated candidates preserves both CPU and CUDA random states and
adds no model evaluation.

The preregistered condition fails. Writing
$G=J_{RR}-J_{CC}=116.866234$, it requires both
$J_{RR}-J_{CR}\geq G/2$ and $J_{RC}-J_{CC}\geq G/2$.
Instead, $J_{CR}$ exceeds $J_{RR}$ by $62.803333$, while the second contrast
recovers $49.9055\%$ of $G$.
In particular, retaining $C$ proposals and restoring only $R$ scoring raises
return from $495.216321$ to $674.885888$.
The factorial interaction
$J_{CC}-J_{CR}-J_{RC}+J_{RR}=-121.125975$ also precludes an additive
allocation of the original loss to the two routes.
This counterexample does not support the proposed claim that damaged
proposal coverage prevents scoring-path recovery in this case.
It does not measure the physical quality of every candidate, demonstrate
improved forecast accuracy, or establish an attack, a defense, or a
population effect. The fixed gate is closed without a seed or threshold sweep.

SWAAP already describes poisoned world-model adaptation and the use of
TD-MPC2 dynamics for both policy-guided proposals and candidate evaluation
\cite{hu2026swaap}. Our distinction is the inference-time factorial separation
of these consumers, not a first demonstration of model poisoning or proposal
dependence. Here that separation contradicts the registered proposal-dominant
explanation. Results and the launch receipt are in
\artifact{runs/tdmpc_proposal_scoring_factorial_v1/summary.json} and
\artifact{runs/tdmpc_proposal_scoring_factorial_v1/launcher.json}; all four
physical traces and the common-first-population captures are retained.
